\section{Right-hand-side perturbations}\label{sec:perturbation}
The coefficient in Theorem~\ref{thm:lower} explodes because the
right-hand side of the relaxed equality lies within $\delta_d$ of the
boundary of the residual image of an improving slice, that is, of the
set of right-hand sides attainable on that slice. This section shows
that a margin from such boundaries controls the coefficient: if the
perturbed problem is feasible and its right-hand side keeps distance more
than $\gamma$ from the boundaries of all residual images, the coefficient
$M/\gamma$ is exact at multiplier zero (Lemma~\ref{lem:rhsmargin}). A right-hand side sampled from a fine
rational grid has such a margin with high probability, which gives a
coefficient of polynomial encoding length
(Theorem~\ref{thm:gridsmoothing}). The principle is classical in
conditioning theory \citep{DunaganSpielmanTeng2011,BuergisserAmelunxen2012}:
a right-hand side far from the boundaries of the convex residual images
is far from any change in the set of feasible integer assignments, and
random perturbations rarely come close to such boundaries. We prove the
cube estimate and its transfer to the grid; the Gaussian variant in
Section~\ref{sec:gaussian} uses \citet[Lemma~2.3.4]{DunaganSpielmanTeng2011}.

Perturbing the right-hand side changes the problem. All guarantees in
this section concern the sampled instance, which may be infeasible and
may have a different optimal integer assignment from the original one.

\subsection{Setting}\label{sec:perturbation-setting}
Let $\mathcal Z$ be a finite index set with $|\mathcal Z|\le K$, where
$K\ge1$. For each $z\in\mathcal Z$, let $X_z$ be a nonempty compact
convex subset of a Euclidean space (the slice), let
$r_z:X_z\to\R^m$ be affine, with $m\ge1$, and let $f_z:X_z\to\R$ be
continuous and convex. Bounds $\underline f\le\overline f$ with
\[
 \underline f\le f_z(x)\le\overline f\qquad(z\in\mathcal Z,\ x\in X_z)
\]
are supplied, and $M=\overline f-\underline f$ denotes the objective range.
The slices, maps and objectives are fixed before any sampling. For a
right-hand side $b\in\R^m$, let
\begin{equation}\label{eq:rhsmodel}
 \begin{aligned}
 v(b)&=\min\{f_z(x):z\in\mathcal Z,\ x\in X_z,\ r_z(x)=b\},\\
 p_\rho(b)&=\min_{z\in\mathcal Z,\,x\in X_z}
             \{f_z(x)+\rho\norm{r_z(x)-b}_\infty\},
 \end{aligned}
\end{equation}
with $v(b)=+\infty$ if no point is feasible. We call $b$ \emph{feasible}
if $v(b)<\infty$, and we write $v_z(b)$ for the minimum of $f_z$ over
$\{x\in X_z:r_z(x)=b\}$ (with $v_z(b)=+\infty$ if this set is empty).

This is the setting of Section~\ref{sec:geometry}: the native set is the
disjoint union of the slices, the objective is $f_z$ on slice $z$, the
residual is $r_z(x)-b$, and the norm is the infinity norm. In particular,
$p_\rho(b)=L_\rho(0)$ for this instance. For a feasible $b$, we write
$\rho_*[b]$ for the least penalty of this instance, with the multiplier
optimized; the square brackets distinguish it from the fixed-multiplier
threshold $\rho_*(\lambda)$. For infeasible $b$ the least penalty is not
defined; where the distribution of $\rho_*[b]$ over random $b$ is
discussed, any nonnegative value may be assigned to infeasible $b$.

For example, in the model of Section~\ref{sec:upper} with convex slices,
each integer assignment $z$ with a nonempty slice gives such a slice,
with $f_z$ the objective and $r_z$ the left-hand side of the linking
equality for this fixed $z$. Then $K$ can be taken as the product of the
numbers of allowed values of the integer variables, so $\log K$ is
polynomial in the input length, but $K$ itself can be exponential in the
number of integer variables. The present setting needs neither a Slater condition nor
quadratic data.

The \emph{residual image} $C_z=r_z(X_z)$ is the set of right-hand sides
attainable on slice $z$; it is nonempty, compact and convex, and the
problem with right-hand side $b$ has a feasible point in slice $z$
exactly when $b\in C_z$.
Boundaries $\partial C$ are topological boundaries in $\R^m$. Thus a
closed convex set with empty interior, such as a singleton, is its own
boundary, and every right-hand side in it has margin zero: arbitrarily
small perturbations leave it. Distances are
$\operatorname{dist}_\infty(b,S)=\inf_{y\in S}\norm{b-y}_\infty$, and
similarly $\operatorname{dist}_2$, with distance $+\infty$ to the empty
set. The latter convention matters only in Lemma~\ref{lem:cubetube},
where $C=\varnothing$ or $C=\R^m$ is allowed.

\subsection{A deterministic boundary-margin bound}
The idea is a convex repair. Suppose every residual image has boundary
distance greater than $\gamma$ from $b$. On a slice whose image contains
$b$, the whole cube $b+[-\gamma,\gamma]^m$ then consists of attainable
right-hand sides; in particular, this image is full-dimensional. A native
point $x$ at residual distance $t=\norm{r_z(x)-b}_\infty>0$ can be
repaired by mixing it with a point $x'$ whose residual lies on the far
side of $b$ at distance $\gamma$. The mixture has residual exactly $b$,
and convexity bounds the cost of the repair by $(M/\gamma)t$. On a slice
whose image does not contain $b$, every residual is at distance more than
$\gamma$ from $b$, so the penalty $(M/\gamma)t$ exceeds the whole
objective range.

\begin{lemma}\label{lem:rhsmargin}
In the setting of Section~\ref{sec:perturbation-setting}, let $b\in\R^m$
be feasible and let $\gamma>0$ satisfy
$\operatorname{dist}_\infty(b,\partial C_z)>\gamma$ for every
$z\in\mathcal Z$. Then $p_{M/\gamma}(b)=v(b)$, that is, the coefficient
$M/\gamma$ is exact at multiplier zero for right-hand side $b$; in
particular $\rho_*[b]\le M/\gamma$. Every coefficient $\rho>M/\gamma$ is
minimizer-exact at multiplier zero.
\end{lemma}
\begin{proof}
We show that $f_z(x)+(M/\gamma)\norm{r_z(x)-b}_\infty\ge v(b)$ for every
slice $z$ and every $x\in X_z$, first on slices whose image contains $b$
and then on the others.

\emph{Slices with $b\in C_z$.} First, $b+[-\gamma,\gamma]^m\subseteq C_z$.
Otherwise some $y$ with $\norm{y-b}_\infty\le\gamma$ lies outside $C_z$.
Since $C_z$ is closed and contains $b$, the last point of $C_z$ on the
segment from $b$ to $y$ lies on $\partial C_z$, at infinity-norm distance
at most $\gamma$ from $b$, a contradiction.

Now let $x\in X_z$ with $t=\norm{r_z(x)-b}_\infty>0$. The point
$b-(\gamma/t)(r_z(x)-b)$ lies in $b+[-\gamma,\gamma]^m\subseteq C_z$, so
some $x'\in X_z$ satisfies
\[
 r_z(x')=b-\frac\gamma t\,(r_z(x)-b).
\]
The convex combination $x_b=(\gamma x+tx')/(\gamma+t)$ lies in $X_z$, and
since $r_z$ is affine, $r_z(x_b)=b$. Convexity of $f_z$ and
$f_z(x')\le\overline f$ give
\[
 v_z(b)\le f_z(x_b)\le\frac\gamma{\gamma+t}f_z(x)
                    +\frac t{\gamma+t}\overline f .
\]
Rearranging, and using $v(b)\le v_z(b)$ and
$\overline f-v_z(b)\le\overline f-\underline f=M$,
\[
 f_z(x)\ge v_z(b)-\frac{\overline f-v_z(b)}{\gamma}\,t
       \ge v(b)-\frac M\gamma\,t .
\]
This inequality also holds at $t=0$, since then $x$ is feasible and
$f_z(x)\ge v(b)$.

\emph{Slices with $b\notin C_z$.} A closest point of the compact set
$C_z$ to $b$ lies on $\partial C_z$, so
$\operatorname{dist}_\infty(b,C_z)=\operatorname{dist}_\infty(b,\partial C_z)>\gamma$.
Hence every $x\in X_z$ has $t=\norm{r_z(x)-b}_\infty>\gamma$ and
\[
 f_z(x)+\frac M\gamma\,t\ge\underline f+M=\overline f\ge v(b).
\]

Thus $p_{M/\gamma}(b)\ge v(b)$, and a feasible optimum attains $v(b)$
with zero penalty, so equality holds. If $\rho>M/\gamma$, every point
with nonzero residual has penalized value strictly greater than $v(b)$,
so the minimizers are exactly the optimal feasible points.
\end{proof}
The argument needs neither differentiability nor computed multipliers.
It is a right-hand-side analogue of Lemma~\ref{lem:slices}: the margin
$\gamma$ replaces the multiplier bound on slices whose image contains
$b$ by the convex repair, and it gives the residual separation $\gamma$
on the other slices. The relation between exact penalties and multiplier bounds in
convex optimization is classical; see
\citet[Theorem~4.2]{FriedlanderTseng2007} for a polar-gauge
characterization under a nonempty compact solution set.

Convexity of the objective is essential. Take one slice $X=[-1,1]$ with
$r(x)=x$, right-hand side $b=0$ and $f(x)=-\sqrt{|x|}$. The image
$[-1,1]$ has boundary distance one from $b$, and the objective range is
one. For $0<t\le1$, the balanced mixture of $x=t$ and $x=-t$ with equal
weights has average penalized value $-\sqrt t+\rho t$, independently of
the multiplier. This is negative for small $t$, so by
Proposition~\ref{prop:mixture}, $D_\rho<0=v$ for every $\rho\ge0$: no
finite penalty is value-exact.

\subsection{Convex-boundary tubes in a cube}
The next lemma bounds the probability that a uniform point of a cube
falls near the boundary of a closed convex set. The bound is linear in
the tube width and does not depend on the set.

\begin{lemma}\label{lem:cubetube}
Let $Q=b_0+[-\sigma,\sigma]^m$ with $\sigma>0$, and let $B$ be uniform on
$Q$. For every closed convex set $C\subseteq\R^m$ and every $\gamma\ge0$,
\begin{equation}\label{eq:cubetube}
 \Pr\{\operatorname{dist}_\infty(B,\partial C)\le\gamma\}
       \le\min\{1,2m\gamma/\sigma\}.
\end{equation}
The same bound holds with $\operatorname{dist}_2$ in place of
$\operatorname{dist}_\infty$.
\end{lemma}
\begin{proof}
The plan is to cover the tube by an outer shell and an inner shell of the
set, bound the volume of each shell inside $Q$ one coordinate direction
at a time, and let the shell width decrease to $\gamma$.

If $C=\varnothing$ or $C=\R^m$, the boundary is empty and the event is
empty. Otherwise $\partial C$ is nonempty. Fix $w>0$ and let
$S_i=[-w,w]e_i$ for $i=1,\ldots,m$. For a closed convex set $E$, the
addition $E+S_i$ and the erosion $E\ominus S_i=\{x:x+S_i\subseteq E\}$
are again closed and convex.

\emph{One-dimensional fact.} Let $I\subseteq\R$ be a closed convex set
(empty, a point, a closed interval, a closed ray or $\R$) and let $W$ be
an interval of length $2\sigma$. Then
\[
 \bigl|((I+[-w,w])\setminus I)\cap W\bigr|\le2w,\qquad
 \bigl|(I\setminus(I\ominus[-w,w]))\cap W\bigr|\le2w,
\]
where $|\cdot|$ is length. Indeed, addition extends $I$ by $w$ at each
finite endpoint, and erosion removes at most $w$ at each finite endpoint
(possibly deleting a short interval entirely); there are at most two
endpoints.

\emph{Shell volumes.} On each line parallel to $e_i$, the section of
$E+S_i$ is the section of $E$ plus $[-w,w]$, the section of $E\ominus S_i$
is the section of $E$ eroded by $[-w,w]$, and the section of $Q$ is an
interval of length $2\sigma$. The one-dimensional fact and Fubini's
theorem over the $(m-1)$-dimensional face of $Q$ orthogonal to $e_i$ give
\begin{align*}
 \operatorname{vol}\bigl(((E+S_i)\setminus E)\cap Q\bigr)
   &\le2w(2\sigma)^{m-1},\\
 \operatorname{vol}\bigl((E\setminus(E\ominus S_i))\cap Q\bigr)
   &\le2w(2\sigma)^{m-1}.
\end{align*}
We measure set differences intersected with $Q$, which have finite volume.
Successive additions of $S_1,\ldots,S_m$ to $C$ give $C+[-w,w]^m$.
Successive erosions give $C\ominus[-w,w]^m$, because
$(E\ominus S)\ominus T=E\ominus(S+T)$. The intermediate sets are closed
and convex, so the bounds apply at each step, and telescoping over the
$m$ steps gives
\[
 \operatorname{vol}\bigl(((C+[-w,w]^m)\setminus C)\cap Q\bigr)\le2mw(2\sigma)^{m-1},
 \quad
 \operatorname{vol}\bigl((C\setminus(C\ominus[-w,w]^m))\cap Q\bigr)\le2mw(2\sigma)^{m-1}.
\]

\emph{Tube containment.} Let $w>\gamma$ and let
$\operatorname{dist}_\infty(x,\partial C)\le\gamma$. Since $\partial C$ is
closed, there is $y\in\partial C$ with $\norm{x-y}_\infty\le\gamma$. As
$y\in C$, we get $x\in C+[-w,w]^m$. As $y$ is a boundary point, there is
an exterior point $u\notin C$ with $\norm{u-y}_\infty<w-\gamma$; then
$\norm{u-x}_\infty<w$, so $x+[-w,w]^m\not\subseteq C$ and
$x\notin C\ominus[-w,w]^m$. Hence the tube lies in the union of the two
shells. Dividing their total volume bound $4mw(2\sigma)^{m-1}$ by
$\operatorname{vol}(Q)=(2\sigma)^m$ gives probability at most $2mw/\sigma$.
Letting $w\downarrow\gamma$ proves \eqref{eq:cubetube}, including
$\gamma=0$. Only closedness and convexity of $C$ were used, so unbounded
and lower-dimensional sets are covered.

Finally, $\norm{\cdot}_\infty\le\norm{\cdot}_2$, so the Euclidean event
is contained in the infinity-norm event.
\end{proof}
The coefficient $2m$ is sharp for bounds of the form $c_m\gamma/\sigma$.
Let $C=b_0+[-a,a]^m$ with $0<a<\sigma$, and let $0<\gamma\le a$ with
$a+\gamma\le\sigma$. The tube is the cube of radius $a+\gamma$ minus the
open cube of radius $a-\gamma$, both centred at $b_0$ and contained in $Q$,
so its probability is $((a+\gamma)^m-(a-\gamma)^m)/\sigma^m$. Its
derivative at $\gamma=0$ is $(2m/\sigma)(a/\sigma)^{m-1}$, which tends to
$2m/\sigma$ as $a\uparrow\sigma$.

\subsection{Rational-grid exactness and its encoding}
To obtain rational right-hand sides of controlled encoding length, we
sample from a finite grid and couple the grid sample with a uniform point
of $Q$. For an integer $q\ge1$, let $h=2\sigma/q$ and sample each
coordinate of $G$ independently and uniformly from the $q$ cell centres
\begin{equation}\label{eq:centergrid}
 b_{0i}-\sigma+(j+1/2)h,\qquad j=0,\ldots,q-1.
\end{equation}
Let $U$ be uniform on $[-h/2,h/2]^m$ and independent of $G$. Then
$B=G+U$ is uniform on $Q$, and $\norm{B-G}_\infty\le h/2=\sigma/q$.
Since $\operatorname{dist}_\infty(\cdot,\partial C)$ is one-Lipschitz in
the infinity norm,
\[
 \{\operatorname{dist}_\infty(G,\partial C)\le\gamma\}\subseteq
 \{\operatorname{dist}_\infty(B,\partial C)\le\gamma+\sigma/q\},
\]
and Lemma~\ref{lem:cubetube} gives, for every closed convex $C$ and
$\gamma\ge0$,
\begin{equation}\label{eq:gridtube}
 \Pr\{\operatorname{dist}_\infty(G,\partial C)\le\gamma\}
       \le\min\{1,2m(\gamma/\sigma+1/q)\}.
\end{equation}
The rounding term $2m/q$ is necessary. Convex boundaries have Lebesgue
measure zero, but grid points can lie on them with positive probability:
for $m=1$, an interval whose two endpoints are distinct grid points has
boundary probability $2/q$ at $\gamma=0$, exactly the correction in
\eqref{eq:gridtube}.

\begin{theorem}\label{thm:gridsmoothing}
Assume the setting of Section~\ref{sec:perturbation-setting}. Let
$b_0\in\Q^m$ and let $\sigma,\epsilon,\underline f,\overline f$ be
rational with $\sigma>0$ and $0<\epsilon<1$. Let $G$ be sampled from the
grid \eqref{eq:centergrid} with
\begin{equation}\label{eq:gridsmoothing}
 q\ge\left\lceil\frac{4mK}{\epsilon}\right\rceil,\qquad
 \gamma=\frac{\sigma\epsilon}{4mK},\qquad
 \bar\rho=\frac M\gamma=\frac{4mKM}{\sigma\epsilon}.
\end{equation}
Then, with probability at least $1-\epsilon$, either the sampled problem
is infeasible or $p_{\bar\rho}(G)=v(G)$. On the same event, if the
sampled problem is feasible, every coefficient $\rho>\bar\rho$ is
minimizer-exact at multiplier zero.
\end{theorem}
\begin{proof}
By \eqref{eq:gridtube} and a union bound over the at most $K$ images,
\[
 \Pr\{\operatorname{dist}_\infty(G,\partial C_z)\le\gamma
      \text{ for some }z\in\mathcal Z\}
 \le2mK\Bigl(\frac\gamma\sigma+\frac1q\Bigr)
 \le2mK\Bigl(\frac\epsilon{4mK}+\frac\epsilon{4mK}\Bigr)=\epsilon.
\]
On the complementary event every boundary distance exceeds $\gamma$, and
Lemma~\ref{lem:rhsmargin} applies whenever $G$ is feasible.
\end{proof}
In words: with probability at least $1-\epsilon$, the grid sample has
boundary distance greater than $\gamma=\sigma\epsilon/(4mK)$ from every
residual image; then either no slice is feasible, or
$\rho_*[G]\le\bar\rho=M/\gamma$.

\begin{remark}[Encoding]\label{rem:gridencoding}
Choose $q$ as the smallest power of two with
$q\ge\lceil4mK/\epsilon\rceil$. Then
$\log_2q=O(1+\log m+\log K+\log(1/\epsilon))$, and exact uniform sampling
uses $m\log_2q$ independent random bits. The grid points and $\bar\rho$
have encoding length polynomial in the supplied rational data
$b_0,\sigma,\epsilon,\underline f,\overline f$, in $m$ and in $\log K$; an
integer coefficient strictly larger than $\bar\rho$, such as
$\lfloor\bar\rho\rfloor+1$, has the same bound, although $\bar\rho$ is
numerically of order $K$. Neither computing $\bar\rho$ nor sampling
requires enumerating the slices. The construction does not check the
structural assumptions of Section~\ref{sec:perturbation-setting} and
does not solve the sampled problem.
\end{remark}

\begin{remark}[Endpoint grid]\label{rem:endpointgrid}
Theorem~\ref{thm:gridsmoothing} also holds for the grid with $q+1$ points
$b_{0i}-\sigma+jh$, $0\le j\le q$, per coordinate. Its jitter cells tile
the cube of radius $\sigma+h/2$, so Lemma~\ref{lem:cubetube} for that cube
gives
\[
 \Pr\{\operatorname{dist}_\infty(G,\partial C)\le\gamma\}
 \le\min\left\{1,\frac{2m(\gamma+h/2)}{\sigma+h/2}\right\}
 \le\min\{1,2m(\gamma/\sigma+1/q)\}.
\]
With $q+1$ points per coordinate, the random-bit count of
Remark~\ref{rem:gridencoding} no longer applies. For either grid the
Euclidean boundary event obeys the same estimate, but it pairs only with
a Euclidean penalty: a Euclidean margin greater than $\gamma$ does not
give an infinity-norm margin greater than $\gamma$.
\end{remark}

\subsection{Feasibility, local stability and the original instance}
\label{sec:perturbation-stability}
\paragraph{The infeasibility alternative.}
Theorem~\ref{thm:gridsmoothing} cannot guarantee feasibility. For
example, if no residual image meets $Q$, every sample is infeasible. If
some slice satisfies $Q\subseteq C_{z_0}$, every grid point is feasible,
and the theorem gives exactness at multiplier zero with probability at
least $1-\epsilon$ unconditionally. More generally, if the grid sample is
feasible with probability at least $p_0>0$, then the conditional failure
probability given feasibility is at most $\epsilon/p_0$.

This bound concerns the grid actually used; refining the grid can lower
its feasibility probability. For a recipe valid for every grid size,
suppose that $B$ uniform on $Q$ is feasible with probability at least a
known rational $p_0>0$. Under the jitter coupling, if $B\in C_z$ and $G\notin C_z$,
then the segment from $G$ to $B$ crosses $\partial C_z$, so
$\operatorname{dist}_\infty(B,\partial C_z)\le\sigma/q$.
Lemma~\ref{lem:cubetube} and a union bound give
\[
 \Pr\{G\text{ feasible}\}\ge p_0-\frac{2mK}q .
\]
Running the construction with $\epsilon p_0/2$ in place of $\epsilon$
(so $q\ge\lceil8mK/(\epsilon p_0)\rceil$ and $\bar\rho=8mKM/(\sigma\epsilon p_0)$)
gives $\Pr\{G\text{ feasible}\}\ge p_0(1-\epsilon/4)$, while the failure
probability is at most $\epsilon p_0/2$. Hence the conditional failure
probability given feasibility is at most
$\epsilon/(2-\epsilon/2)<\epsilon$.

\paragraph{Local stability.}
Let $b$ be feasible with
$\operatorname{dist}_\infty(b,\partial C_z)>\gamma$ for every $z$. By the
first step of the proof of Lemma~\ref{lem:rhsmargin}, every residual
image $C_z$ that contains $b$ contains the whole cube
$b+[-\gamma,\gamma]^m$, while every other image has distance greater than
$\gamma$ from $b$. Hence the set of slices (integer assignments) that are
feasible for the right-hand side is the same at every point of the open cube
of radius $\gamma$ about $b$. We use the smaller cube
$b+[-\gamma/2,\gamma/2]^m$ because, by the one-Lipschitz property of the
distance, every point $b'$ of it keeps margin greater than $\gamma/2$.
Applying the repair inequality of Lemma~\ref{lem:rhsmargin} at $b'$ with
margin $\gamma/2$ to an optimal point for $b''$, and then with the roles
exchanged, gives
\[
 |v_z(b')-v_z(b'')|\le\frac{2M}\gamma\norm{b'-b''}_\infty
\]
for every slice $z$ whose image contains $b$ and all $b',b''$ in the
cube. A minimum over
the same finite collection of slice values keeps this Lipschitz bound, so
$v$ is $2M/\gamma$-Lipschitz in the infinity norm on the cube. By
Lemma~\ref{lem:rhsmargin} with margin $\gamma/2$, the coefficient
$2M/\gamma$ is exact at multiplier zero at every point of the cube, and
every larger coefficient is minimizer-exact at multiplier zero there.
Lipschitz value bounds across discrete choices under uniform Slater
conditions also appear in mixed-integer optimal control
\citep[Theorem~3]{GugatHante2017}.

\paragraph{The original instance.}
The local conclusion does not recover the optimum at $b_0$. The chain of
Section~\ref{sec:lower} fits the present setting with $m=1$, two slices
and $M=1$: the slice with binary variable zero has residual image
$[-1,1]$ and cost $0$, and the improving slice (binary variable one) has
residual image $[\delta_d,1]$ and cost $-1$. The right-hand side $0$ lies
within $\delta_d$ of the boundary of $[\delta_d,1]$, so
Lemma~\ref{lem:rhsmargin} gives only coefficients above $1/\delta_d$,
the zero-multiplier threshold of Section~\ref{sec:lower}; this is why the
coefficient explodes. Moving the right-hand side to $\delta_d$ makes the
improving assignment feasible, and the optimum becomes $-1$ instead of
$0$. Short penalties after perturbation can therefore come with a
different optimal integer assignment.

\subsection{Numerical size: unavoidable tails and slice dependence}
\label{sec:numerical-size}
The coefficient $\bar\rho=4mKM/(\sigma\epsilon)$ of
Theorem~\ref{thm:gridsmoothing} has short encoding but can be
numerically large. The following examples show which parts of this size
are unavoidable:
\begin{itemize}
 \item[(a)] the order $K/\epsilon$ is necessary: some instances satisfy
 $\Pr\{\rho_*[B]>T\}\ge(K-1)/T$ for $T\ge K$
 (Example~\ref{ex:tail});
 \item[(b)] a lower bound of order $K$ can hold almost surely, not just
 with small probability (Example~\ref{ex:slicecount});
 \item[(c)] two slices can give $\mathbb E\rho_*[B]=+\infty$ while
 $\mathbb E\log(1+\rho_*[B])$ is finite (Example~\ref{ex:infinitemean});
 \item[(d)] an exceptional grid point can have infinite threshold, so
 the high-probability guarantee cannot be turned into a bound on
 expectations (Example~\ref{ex:atom}).
\end{itemize}
In all examples $m=1$, so the norm is the absolute value, and
$\underline f=-1$, $\overline f=0$, $M=1$. All thresholds are optimized
thresholds $\rho_*[b]$ unless we say otherwise; since exactness at
multiplier zero implies value-exactness, lower bounds on $\rho_*[b]$ also
bound the zero-multiplier threshold. The examples list their slices
explicitly, as the slice model allows; they are not derived from a
succinct mixed-integer encoding.

The lower bounds use one balanced mixture
(Proposition~\ref{prop:mixture}). Suppose a slice with cost $-1$ has a
singleton residual image at distance $t>0$ from $b$, and a cost-zero
slice contains a point whose residual $r_z(x)-b$ has the opposite sign
and magnitude $s>0$. Put weight
$s/(s+t)$ on the singleton point and weight $t/(s+t)$ on the cost-zero
point. The residuals cancel, and the average penalized value is
\begin{equation}\label{eq:tailmixture}
 \frac{s}{s+t}(-1+2\rho t)
\end{equation}
for every multiplier. If $v(b)=0$, this mixture shows
$\rho_*[b]\ge1/(2t)$.

\begin{example}[The order $K/\epsilon$ is necessary]\label{ex:tail}
Let an anchor slice have residual image $[0,1]$ and cost zero, and let
$K-1$ singleton slices, $K\ge2$, have images $\{j/K\}$,
$1\le j\le K-1$, and cost $-1$. At every right-hand side
$b\in(0,1)$ that is not one of the points $j/K$, the primal value is
zero, and \eqref{eq:tailmixture} with the nearest singleton at distance
$t$ gives $\rho_*[b]\ge1/(2t)$. Let $B$ be uniform on $[0,1]$. For
$T\ge K$, the intervals $|B-j/K|<1/(2T)$ are disjoint, lie in $(0,1)$ and
have total length $(K-1)/T$. Consequently
\begin{equation}\label{eq:penaltytail}
 \Pr\{\rho_*[B]>T\}\ge\frac{K-1}T\qquad(T\ge K).
\end{equation}
Let $0<\epsilon<(K-1)/K$. For $K\le T<(K-1)/\epsilon$,
\eqref{eq:penaltytail} gives $\Pr\{\rho_*[B]>T\}>\epsilon$, and for
$T<K$ the same holds by monotonicity. Hence the $(1-\epsilon)$-quantile
of $\rho_*[B]$ is at least $(K-1)/\epsilon$. Here $b_0=1/2$,
$\sigma=1/2$ and $M=1$, so Theorem~\ref{thm:gridsmoothing} gives the
coefficient $8K/\epsilon$; the order $K/\epsilon$ cannot be improved, even
though every sample is feasible and the multiplier is optimized.

The grid sample of Theorem~\ref{thm:gridsmoothing} behaves the same way.
On $[0,1]$ the cell centres are $(i+1/2)/q$; let $K$ divide $q$, so that
no centre is a singleton image. If $q\ge2K$, each singleton image $j/K$
has two adjacent centres at distance $1/(2q)$, and these $2(K-1)$ centres
are distinct; hence $\Pr\{\rho_*[G]\ge q\}\ge2(K-1)/q$. More generally,
for $K\le T\le q$, each singleton image has at least $q/T-1$ centres at
distance less than $1/(2T)$, these sets are disjoint for different $j$,
and at each such centre $\rho_*[G]>T$. Hence
$\Pr\{\rho_*[G]>T\}\ge(K-1)(1/T-1/q)$. If $\epsilon\le(K-1)/(2K)$ and
$q\ge4K/\epsilon$, then $T=(K-1)/(2\epsilon)$ satisfies $K\le T\le q$, and
$\Pr\{\rho_*[G]>T\}\ge2\epsilon-(K-1)/q>\epsilon$. So the
$(1-\epsilon)$-quantile of $\rho_*[G]$ is at least $(K-1)/(2\epsilon)$,
and the order $K/\epsilon$ is necessary for the grid result as well.
\end{example}

\begin{example}[An almost-sure slice-count obstruction]\label{ex:slicecount}
Keep the anchor slice with image $[0,1]$ and cost zero, and use $J+1$
singleton slices, $J\ge1$, with images $\{j/J\}$, $j=0,\ldots,J$, and
cost $-1$. There are $K=J+2$ slices. Let $b\in(0,1)$ lie strictly between
the neighbouring singleton images $j/J$ and $(j+1)/J$, and let
$g_1=b-j/J>0$ and $g_2=(j+1)/J-b>0$ be its distances to them, so
$g_1+g_2=1/J$. The primal value is zero. Put weight $g_2/(g_1+g_2)$ on
the lower singleton and $g_1/(g_1+g_2)$ on the upper one. The residuals
cancel, the cost is $-1$, and the mean absolute residual is
\[
 \frac{2g_1g_2}{g_1+g_2}=2Jg_1g_2\le\frac1{2J}.
\]
Hence $\rho_*[b]\ge2J=2(K-2)$ for every $b\in(0,1)$ other than the
singleton images, that is, for almost every $b\in[0,1]$.
\end{example}

\begin{example}[Infinite mean with two slices]\label{ex:infinitemean}
Take an anchor slice with image $[-1,1]$ and cost zero and a singleton
slice with image $\{0\}$ and cost $-1$. For $b\in(-1,1)\setminus\{0\}$,
the primal value is zero, and \eqref{eq:tailmixture} with $t=|b|$ gives
$\rho_*[b]\ge1/(2|b|)$. Conversely, at $\rho=1/(2|b|)$ and multiplier
$\lambda=-\rho\operatorname{sign}(b)$, the singleton's penalized value is
$-1+\lambda(-b)+\rho|b|=-1+2\rho|b|=0$, and the anchor's penalized value
$\lambda e+\rho|e|$ at residual $e$ is nonnegative. Hence
\begin{equation}\label{eq:infinitemean}
 \rho_*[b]=\frac1{2|b|}.
\end{equation}
The zero-multiplier threshold is instead $1/|b|$. For $B$ uniform on
$[-1,1]$, we get $\mathbb E\rho_*[B]=+\infty$, while
$\mathbb E\log(1+\rho_*[B])<\infty$ because $\log(1+1/(2|b|))$ is
integrable near zero.
\end{example}

\begin{example}[An exceptional grid point with infinite threshold]
\label{ex:atom}
Take a convex quadratic slice
$\{(\alpha,\beta):0\le\alpha\le1,\ \alpha^2\le\beta\le1\}$ with cost
$-\alpha$ and residual $\beta$, so its image is $[0,1]$, and a singleton
slice with image $\{-1\}$ and cost zero. At right-hand side $0$ the
constraint $\beta=0$ forces $\alpha=0$, so the primal value is zero. For
$0<u\le1$, mix the point $(\alpha,\beta)=(u,u^2)$ with weight
$1/(1+u^2)$ and the singleton with weight $u^2/(1+u^2)$. The residuals
cancel, the cost is $-u/(1+u^2)$, and the mean absolute residual is
$2u^2/(1+u^2)$. By Corollary~\ref{cor:ratio}, $\rho_*[0]$ is at least the
ratio $1/(2u)$ for every such $u$, so $\rho_*[0]=+\infty$: no finite
penalty is value-exact. A rational grid containing $0$, for instance the
grid \eqref{eq:centergrid} with $b_0=0$ and odd $q$, gives this
right-hand side positive probability. Hence $\mathbb E\rho_*[G]=+\infty$
(with any nonnegative value assigned to infeasible samples), and no
almost-sure bound on the threshold or its encoding holds. The margin
condition of Lemma~\ref{lem:rhsmargin} excludes this point, since $0$
lies on the boundary of $[0,1]$. The refined Slater condition also fails
there: on $\beta=0$ the inequality $\alpha^2\le\beta$ cannot hold strictly.
\end{example}

\subsection{A Gaussian companion in the Euclidean norm}\label{sec:gaussian}
Lemma~\ref{lem:rhsmargin} and its proof hold for any norm on $\R^m$,
provided the same norm defines both the margin and the penalty: replace
the cube $b+[-\gamma,\gamma]^m$ by the ball of radius $\gamma$. For
Gaussian noise we use the Euclidean norm for both. There is no grid
here, so this result has no encoding statement.

\begin{proposition}\label{prop:gaussian}
Assume the setting of Section~\ref{sec:perturbation-setting}, but with
the Euclidean norm in the penalty of \eqref{eq:rhsmodel}. Let
$B=b_0+\sigma\xi$ with $\xi\sim\mathcal N(0,I_m)$, $\sigma>0$ and
$0<\epsilon<1$. Then for every $\gamma\ge0$,
\begin{equation}\label{eq:gaussiantube}
 \Pr\{\min_{z\in\mathcal Z}\operatorname{dist}_2(B,\partial C_z)\le\gamma\}
       \le\frac{8Km^{1/4}\gamma}\sigma .
\end{equation}
Consequently, with probability at least $1-\epsilon$, either the sampled
problem is infeasible or the coefficient
\begin{equation}\label{eq:gaussianpenalty}
 \frac{8MKm^{1/4}}{\sigma\epsilon}
\end{equation}
is exact at multiplier zero for the Euclidean penalty; every larger
coefficient is then minimizer-exact at multiplier zero.
\end{proposition}
\begin{proof}
First let $C$ be a nonempty compact convex set with nonempty interior.
\citet[Lemma~2.3.4]{DunaganSpielmanTeng2011} bound the probabilities of
the outer tube $\{\operatorname{dist}_2(B,\partial C)\le\gamma\}\setminus C$
and of the inner tube $\{\operatorname{dist}_2(B,\partial C)\le\gamma\}\cap C$
each by $4m^{1/4}\gamma/\sigma$. Hence
$\Pr\{\operatorname{dist}_2(B,\partial C)\le\gamma\}\le8m^{1/4}\gamma/\sigma$.

Now let $C$ be nonempty, compact and convex with empty interior, so
$\partial C=C$. For $\eta>0$, the parallel body
$C_\eta=C+\eta\mathbb B_2$ is compact and convex with nonempty interior.
We claim $\operatorname{dist}_2(x,\partial C_\eta)\le|\operatorname{dist}_2(x,C)-\eta|$.
Let $p$ be the projection of $x$ onto $C$, and let $u$ be a unit vector
in the normal cone of $C$ at $p$: take $u=(x-p)/\norm{x-p}_2$ if
$x\notin C$, and a unit normal to the affine hull of $C$ if $x\in C$
(it exists because $C$ has empty interior). Then $p+su$ has projection
$p$ and distance $s$ from $C$ for every $s\ge0$. So $p+\eta u\in C_\eta$,
while $p+su\notin C_\eta$ for $s>\eta$; thus $p+\eta u\in\partial C_\eta$.
Since $x=p+\operatorname{dist}_2(x,C)\,u$, the claim follows. Hence
$\operatorname{dist}_2(x,C)\le\gamma$ implies
$\operatorname{dist}_2(x,\partial C_\eta)\le\gamma+\eta$, and the
first case gives
\[
 \Pr\{\operatorname{dist}_2(B,\partial C)\le\gamma\}
 \le\Pr\{\operatorname{dist}_2(B,\partial C_\eta)\le\gamma+\eta\}
 \le\frac{8m^{1/4}(\gamma+\eta)}\sigma .
\]
Letting $\eta\downarrow0$ gives the same bound as for convex bodies.

A union bound over the at most $K$ images proves \eqref{eq:gaussiantube}.
With $\gamma=\sigma\epsilon/(8Km^{1/4})$, the right-hand side equals
$\epsilon$. On the complementary event, every Euclidean boundary distance
exceeds $\gamma$, and the Euclidean version of Lemma~\ref{lem:rhsmargin}
shows that $M/\gamma$, which equals \eqref{eq:gaussianpenalty}, is exact
at multiplier zero whenever $B$ is feasible, and that every larger
coefficient is minimizer-exact at multiplier zero.
\end{proof}
